\chapter{Взаємодія плазми зі стінкою і дивертор}\label{ch:wall}

За сепаратрисою силові лінії закінчуються на мішенях дивертора. Контакт плазми з твердим тілом визначає теплове навантаження на мішені, джерело домішок для ядра (розд.~\ref{ch:transport}) і граничні умови для потоків SOL (п.~\ref{sec:06-sol}).

\begin{corpusexample}[лімітер і перша стінка JET: проєкт 1975 і машина 1983]
Перші великі токамаки працювали без дивертора: межу плазми задавав \emph{лімітер} (limiter) --- тіло, якого першим торкається плазма. Проєкт JET 1975~р. передбачав три осесиметричні «рейкові» лімітери з пластин, положення яких можна регулювати ззовні без розгерметизації, з газовим охолодженням~\cite[с.~301]{JET1975}. Зовнішню рейку мали скласти 16 молібденових пластин $100\times60$~см завтовшки \SI{1}{см} (разом близько \SI{1}{т} Mo). Розмір пластини обрано з простої оцінки: якщо близько \SI{10}{МДж}, які плазма втрачає за першу секунду розряду, рівномірно вкласти в одну таку пластину, її поверхня нагріється приблизно до \SI{1000}{\celsius}~\cite[с.~314]{JET1975}. Камеру з інконелю для досягнення базового тиску мали прогрівати до \SI{500}{\celsius} гарячим інертним газом у міжстінному просторі~\cite[с.~299, 301]{JET1975}. На машині, що запрацювала 1983~р., рейку замінили 12 дискретними модулями по \SI{0.32}{м^2} на зовнішньому екваторі камери: 4 неохолоджувані графітові і 8 водоохолоджуваних з нікелевим покриттям \SI{1.5}{мм}, розрахованих на \SI{e7}{Вт/м^2}; прогрів камери \SI{500}{\celsius} і чистка тліючим розрядом при \SI{250}{\celsius} (п.~\ref{sec:08-gdc}) лишились штатними режимами~\cite[с.~24]{JETJU1983}. Графітові модулі мали розміри $40\times80$~см, поверхню, злегка вигнуту в тороїдальному напрямку для розподілу теплового навантаження, і нагрівались у розрядах приблизно до \SI{800}{\celsius}~\cite[с.~135]{Wesson1999}. Поява графіту поряд із металом (Mo у проєкті, Ni у 1983~р.) --- той самий компроміс «теплове навантаження проти домішок», що розібрано в п.~\ref{sec:08-erosion} (детально --- розд.~\ref{ch:jet}).
\end{corpusexample}

\section{Шар, передшар і перенесення тепла вздовж поля}\label{sec:08-sheath}

Рухливіші електрони заряджають стінку від'ємно, доки струм на плаваючу поверхню не стане нульовим. Квазінейтральність порушується лише в тонкому \emph{шарі Дебая} (Debye sheath) товщиною $\sim\lambda_D=\sqrt{\varepsilon_0T_e/(n_ee^2)}\ll\SI{0,1}{мм}$~\cite[с.~23]{AhoMantila2011}. Лінії падають на мішень під кутом 86--90$^\circ$ до нормалі, тож перед шаром є квазінейтральний \emph{магнітний передшар} (MPS) завтовшки кілька іонних ларморових радіусів ($\sim\SI{1}{мм}$), де й відбувається більша частина падіння потенціалу~\cite[с.~23]{AhoMantila2011}.

\begin{derivation}[критерій Бома і плаваючий потенціал]
\emph{Бом.} Холодні іони входять у шар зі швидкістю $u_0$, електрони больцманівські: $n_i=n_0u_0/u$, $\tfrac12m_iu^2=\tfrac12m_iu_0^2-e\Phi$, $n_e=n_0\ee^{e\Phi/T_e}$. Розвинемо права частину рівняння Пуассона $\varepsilon_0\Phi''=e(n_e-n_i)$ за малим $\Phi<0$ біля межі шару: $e(n_e-n_i)\approx n_0e^2\Phi\,(1/T_e-1/m_iu_0^2)$. Монотонний спад потенціалу можливий, лише якщо коефіцієнт при $\Phi$ невід'ємний (інакше розв'язок осцилює), тобто $u_0^2\ge T_e/m_i$. Для теплих іонів маємо $\mathcal{M}=u/c_s\ge1$, $c_s=\sqrt{(ZT_e+\gamma_{\mathrm a}T_i)/m_i}$, де $\gamma_{\mathrm a}=1$--3 залежно від колізійності й геометрії~\cite[с.~24]{AhoMantila2011}.
\emph{Передшар.} Щоб прискорити іони від $\mathcal{M}\approx0$ до $\mathcal{M}=1$, потрібне слабке поле вздовж лінії. Ізотермічний закон збереження імпульсу $n(1+\mathcal{M}^2)=\mathrm{const}$ дає $n_{se}=n_0/2$, звідки $e\Delta\Phi_{\mathrm{ps}}=T_e\ln\tfrac12\approx-0{,}7\,T_e$, як і в~\cite[с.~24]{AhoMantila2011}.
\emph{Плаваюча стінка.} Прирівняємо іонний потік $\Gamma_i=n_{se}c_s$ до електронного $\tfrac14n_{se}\bar v_e\,\ee^{eV_{\mathrm{sh}}/T_e}$, $\bar v_e=\sqrt{8T_e/\pi m_e}$. При $Z=\gamma_{\mathrm a}=1$
\begin{equation}\label{eq:08-Vsh}
  \frac{eV_{\mathrm{sh}}}{T_e}=\ln\frac{4c_s}{\bar v_e}=0{,}5\ln\Big[2\pi\frac{m_e}{m_i}\Big(1+\frac{T_i}{T_e}\Big)\Big],
\end{equation}
що точно збігається з~\cite[с.~23, (2.1)]{AhoMantila2011} (без вторинної емісії і відбиття).
\end{derivation}

Для D при $T_i=T_e$ з \eqref{eq:08-Vsh} маємо $eV_{\mathrm{sh}}=-2{,}84\,T_e$, при $T_i=0$~--- $-3{,}19\,T_e$. Повне падіння $\approx3{,}5\,T_e/e$ задає $\Er$ у SOL (п.~\ref{sec:06-sol}). Для стінки важливіше, що шар розганяє багатозарядні важкі іони до енергій саморозпилення (п.~\ref{sec:08-erosion}).

\textbf{Передача тепла.} Кожна пара іон--електрон, що входить у шар з потоком $\Gamma_{se}$, виносить на мішень енергію $7$--$8\,T_e$~\cite[с.~25]{AhoMantila2011}. Бар'єр $e|V_{\mathrm{sh}}|$ долають лише найшвидші електрони, тому шар діє як фільтр: розподіл у SOL втрачає високоенергетичний хвіст, а енергію, яку електрони віддають полю шару, отримують прискорені іони. У підсумку електрони SOL охолоджуються~\cite[с.~25]{AhoMantila2011}.

\begin{outofcorpus}[енергія удару, коефіцієнт передачі тепла, теплопровідність Спітцера]
Іон заряду $Z$ набирає в шарі й передшарі $\approx3ZT_e$ і приносить теплову енергію $\approx2T_i$, тож стандартна оцінка (Стенгбі) енергії удару $E_{\mathrm{imp}}\approx3ZT_e+2T_i$ (з точним падінням $3{,}53\,T_e$ для D з \eqref{eq:08-Vsh} на 10--15\,\% більше). $Q_{\parallel t}=\gamma_{\mathrm{sh}}\,n_{se}T_ec_s$, де $\gamma_{\mathrm{sh}}\approx2{,}5\,T_i/T_e+2/(1-\delta_e)-eV_{\mathrm{sh}}/T_e$. Тут $2{,}5\,T_i$~--- енергія іона на вході в шар, $2T_e$~--- напівмаксвеллівський потік енергії електронів, $e|V_{\mathrm{sh}}|$~--- енергія, яку електрони втрачають у шарі і яку отримують іони; $\delta_e$~--- коефіцієнт вторинної емісії. Для D при $T_i=T_e$, $\delta_e=0$ маємо $\gamma_{\mathrm{sh}}=2{,}5+2+2{,}84=7{,}3$. Це узгоджується з <<7--8>> корпусу. Спітцер: $Q_{\parallel e}=-k_{0e}T_e^{5/2}\,\dd T_e/\dd s$, $k_{0e}\approx2000$ (Вт, м, еВ), без залежності від $n_e$.
\end{outofcorpus}

У корпусі $Q_{\parallel s}=-\chis{\parallel s}\grad_\parallel T_s$, $\chis{\parallel e}\propto T_e^{5/2}/n_e\propto\lambda_{\mathrm{coll},e}v_{\mathrm{th},e}$~\cite[с.~26, (2.3)--(2.4)]{AhoMantila2011}; це дифузійність, а провідність $n_e\chis{\parallel e}\propto T_e^{5/2}$. При $\lambda_{\mathrm{coll},e}\to\infty$ формула дає $|Q_{\parallel e}|\to\infty$~--- рідинний опис непридатний~\cite[с.~27]{AhoMantila2011}.

\begin{outofcorpus}[двоточкова модель SOL]
Розглянемо силову трубку довжиною $L$ від середньої площини (\emph{u}) до мішені (\emph{t}). Об'ємних втрат немає, $Q_\parallel$ стале, а тепло переноситься теплопровідністю. Інтегрування закону Спітцера дає $T_u^{7/2}=T_t^{7/2}+\tfrac72Q_\parallel L/k_{0e}$. Разом зі збереженням тиску $n_uT_u=2n_tT_t$ ($\mathcal{M}=1$ на мішені) і умовою шару $Q_\parallel=\gamma_{\mathrm{sh}}n_tT_tc_s(T_t)$, де $c_s(T_t)=\sqrt{2eT_t/m_i}$ ($T_i=T_e$, $T$ в еВ), це замкнена система. При $T_t\ll T_u$ температура $T_u\approx(7Q_\parallel L/2k_{0e})^{2/7}$ майже не залежить від густини, і ($T$ в еВ)
\begin{equation}\label{eq:08-2pm}
  T_t=\frac{m_i}{2e}\Big(\frac{2Q_\parallel}{\gamma_{\mathrm{sh}}\,e\,n_uT_u}\Big)^2\propto n_u^{-2},\qquad n_t=\frac{n_uT_u}{2T_t}\propto n_u^{3}
\end{equation}
\end{outofcorpus}

Скейлінг $n_t\propto n_u^3$ режиму високого рециклінгу є в корпусі; у режимі низького рециклінгу (обмеженому шаром) $n_t\propto n_u$, $\grad_\parallel T\approx0$, $T_t>\SI{20}{еВ}$~\cite[с.~28]{AhoMantila2011}.

\begin{corpusexample}[кінетика електронів, ASCOT на тлі SOLPS5.0 для ASDEX Upgrade~\cite{AhoMantila2011}]
Електрони із зовнішньої середньої площини вели до мішеней крізь зіткнення й поле $\Epar$. На внутрішній трубці з великим $\grad_\parallel T_e$ під X-точкою ($T_e<\SI{5}{еВ}$ на мішенях) надтеплові електрони несли \textbf{70\,\%} енергії, що надходила на зовнішню мішень; частина з них безколізійно приходить з області X-точки~\cite[с.~58]{AhoMantila2011}. Усереднена вздовж лінії колізійність на всіх трубках була однакова, тож критерієм кінетичних ефектів є локальний $\grad_\parallel T_e$, а не вона.
\end{corpusexample}

\section{Рециклінг, нейтрали, детачмент}\label{sec:08-recycling}

Іони рекомбінують на мішенях і повертаються нейтралами, тож джерело частинок SOL зосереджене біля мішеней. Ознака \emph{детачменту} (detachment)~--- «перегин»: густина розряду росте, а потік на мішень, щойно $T_t$ падає до кількох еВ, \emph{спадає}. Для цього повний тиск уздовж лінії не має зберігатися; його знімають об'ємна рекомбінація, зіткнення іонів з нейтралами й випромінювання~\cite[с.~29]{AhoMantila2011}. Поріг моделі кількісно не відтворюють; для ITER закладено частковий детачмент, бо повний шкодить утриманню й загрожує зривом~\cite[с.~29]{AhoMantila2011}. Ненасичені стінки відкачують частинки; після насичення рециклінг~--- 100\,\%~\cite[с.~27--28]{AhoMantila2011}.

\textbf{SOLPS5.0 = B2.5 + Eirene.} B2.5 розв'язує модифіковані рівняння Брагінського для електронів, іонів і домішок на квазіортогональній 2D-сітці CARRE, з дрейфами $\ExB$ і діамагнітним~\cite[с.~48--51]{AhoMantila2011}:
\begin{equation}\label{eq:08-b25}
  \frac{\pa n}{\pa t}+\grad\cdot(n\uvec)=S_n,\qquad
  \frac{\pa(m_in\uvec)}{\pa t}+\grad\cdot(m_in\uvec\uvec)=-\grad p-\grad\cdot\visc_i+\vect{J}\times\Bvec+S_{mu},
\end{equation}
і рівняння енергії для $T_i$ та $T_e$ з обміном енергією між ними. Тут $\visc_i$~--- паралельний тензор в'язкості (у джерелі $-\grad(p+\Pi_i)$)~\cite[с.~48--49, (4.2)--(4.5)]{AhoMantila2011}. Радіальні коефіцієнти аномальні й підганяються під профілі. Eirene розв'язує 3D кінетичне рівняння для нейтралів (Монте-Карло) й ітераційно зшивається з B2.5; на вході в MPS стоять умови Бома--Шодури~\cite[с.~50--51]{AhoMantila2011}. Надійні розв'язки для ASDEX Upgrade є лише для низького рециклінгу з дрейфами (п.~\ref{sec:06-sol}).

\begin{corpusexample}[0D-модель SOL/дивертора у тренажері Fenix~\cite{Fable2022}]
У Fenix є 0D-баланс частинок SOL/дивертора (джерело~--- напуск газу, стік~--- відкачка) і 1D-рівняння розсіяння тепла вздовж трубки від середньої площини до мішені~\cite[с.~4]{Fable2022}. Для ASDEX Upgrade густина на сепаратрисі зростає приблизно як $(\text{напуск})^{1/3}$, густина нейтралів на сепаратрисі $\sim(2\text{--}3)\cdot10^{15}$~м$^{-3}$, а відношення густин дивертор/SOL $\sim30$~\cite[с.~7]{Fable2022}. У EU-DEMO пропущена пелета посилає в SOL хвилю частинок і тепла; для утримання детачменту постійна (feedforward) подача Ar краща за зворотний зв'язок, бо гладкого надійного параметра стану «прикріплений/відірваний» поки немає~\cite[с.~8]{Fable2022}.
\end{corpusexample}
\emph{Наша інтерпретація} (не з корпусу): якщо відкачка пропорційна густині в диверторі, $\Gamma_{\mathrm{puff}}\propto n_t\propto n_u^3$ з \eqref{eq:08-2pm}. Звідси $n_u\propto\Gamma_{\mathrm{puff}}^{1/3}$, що узгоджується зі скейлінгом Fenix.

\section{Ерозія і міграція домішок}\label{sec:08-erosion}

\textbf{Фізичне розпилення} (physical sputtering) (коефіцієнт $Y_{\mathrm{phys}}$) потребує передачі атому енергії понад $E_B$ (\SI{3,38}{еВ} для Be, \SI{8,8}{еВ} для W); через каскад зіткнень поріг $E_{\mathrm{th}}\gg E_B$. Водню для W потрібні сотні еВ, для C~--- менше \SI{50}{еВ}. Тому W ерозують здебільшого нейтрали перезарядки, домішки та саморозпилення багатозарядними іонами W, прискореними шаром~\cite[с.~32]{AhoMantila2011}. Оцінка з п.~\ref{sec:08-sheath}: при $T_e=T_i=\SI{20}{еВ}$ D$^+$ має $E_{\mathrm{imp}}\approx\SI{100}{еВ}$ (нижче порогу для W), а W$^{4+}$~--- вже $\approx\SI{280}{еВ}$. \textbf{Хімічна ерозія} C (вуглеводні з $E_B\sim\SI{1}{еВ}$ проти \SI{7,4}{еВ}) дає $Y=0{,}5\,\%$ уже при енергіях нижче \SI{1}{еВ} (фізичне~--- понад \SI{50}{еВ}); $Y_{\mathrm{chem}}$ максимальний при 700--900~К і спадає з густиною потоку як $\Gamma^{-0{,}54}$, що на порядок зменшило прогноз ерозії мішеней ITER~\cite[с.~33--34]{AhoMantila2011}.

\textbf{Міграція $^{13}$C} (ASDEX Upgrade, ERO на тлі SOLPS5.0). Більшість вуглеводнів не виходить за MPS і осідає біля отвору інжекції; «хвіст» осадження утворюють частинки, що вийшли за MPS і повністю дисоціювали до іонів C~\cite[с.~64]{AhoMantila2011}. Радіальне й полоїдальне поля $\vect E$ у диверторі дають дрейф $\ExB$ уздовж мішені й по нормалі до неї; у прямому полі він збільшує локальне осадження і осадження вище за потоком. За однакового фону плазми зміна лише напрямку $\Btor$ (за будь-якого знака $\Ip$) змінює модельну ефективність осадження з 23--24\,\% до 12--13\,\%~\cite[с.~63--64, рис.~5.2]{AhoMantila2011}. Виміряні значення (24--32\,\% і вдвічі менше в оберненому полі) наведено в п.~\ref{sec:06-sol}.

\begin{corpusexample}[вольфрамові перші дзеркала ITER~\cite{Belyaeva2016}]
Енергії атомів перезарядки, що бомбардують діагностичні дзеркала ($10$--$10^3$~еВ), здебільшого на порядок вищі за енергію зв'язку в ґратці, тож дзеркала неминуче розпилюються~\cite[с.~3]{Belyaeva2016}. Імітація: нейтрони~--- іони W$^{6+}$ 20~МеВ (3~зна), атоми перезарядки~--- Ar 600~еВ~\cite[с.~5]{Belyaeva2016}. Після зняття $\approx\SI{4}{мкм}$ коефіцієнт відбиття рекристалізованого W (W-rc) змінився лише в межах 2\,\%, а в W ITER-grade (W-IG) погіршився на $\approx20\,\%$ через розвинену шорсткість~\cite[с.~23--24]{Belyaeva2016}. Швидкість розпилення не залежить від типу W і опромінення, але залежить від орієнтації зерна: [111] розпилюється найповільніше, [110]~--- найшвидше, бо вздовж [111] ґратка найменш щільно «закриває» шлях пучка й іони заходять глибше~\cite[с.~24]{Belyaeva2016}.
\end{corpusexample}

W, що дістався до педесталу, далі підлягає неокласичному пінчу і перенесенню ELM (п.~\ref{sec:05-imp}): джерело на стінці і вміст W у ядрі пов'язує ланцюжок «шар $\to$ SOL $\to$ педестал».

\section{Пил}\label{sec:08-dust}

Відшаровані осади дають пил~--- мобілізовуваний і з тритієм~\cite[с.~35]{AhoMantila2011}. DUSTT трасує пилинки (Монте-Карло, 3D) на тлі плазми UEDGE під дією сили Лоренца, тертя з іонами й нейтралами та гравітації; радіус $R_d$, температура $T_d$ і потенціал еволюціонують за~\cite[с.~4--5]{Pigarov2006}
\begin{equation}\label{eq:08-dust}
  N_d\frac{\dd R_d}{\dd t}=\Gamma_{\mathrm{in}}-\Gamma_{\mathrm{out}},\qquad
  \frac{\dd}{\dd t}(C_dM_dT_d)=4\pi R_d^2(Q_{\mathrm{in}}-Q_{\mathrm{out}}),\qquad
  J_e-J_{\mathrm{see}}-J_{\mathrm{the}}=J_i ,
\end{equation}
де $N_d$~--- атомна густина речовини (у джерелі з $\rho_d$, $m_d$); $\Gamma_{\mathrm{out}}$~--- відбиття, розпилення і сублімація; потенціал задає баланс струмів з вторинною й термоемісією, падаючі потоки~--- теорія OML~\cite[с.~5]{Pigarov2006}. Ромащенко (дугова технологічна плазма) використовує ту саму схему: баланс струмів $I_i+I_e+I_b(1-\delta)=0$, ємнісний заряд $Z_de=4\pi\varepsilon_0a(1+a/\lambda_D)\Phi_d$~\cite[с.~10, (2), (4)]{Romashchenko2021}, енергобаланс нагріву частинками проти охолодження випромінюванням, емісією і випаровуванням~\cite[с.~13, (9)--(10)]{Romashchenko2021}; OML застосовна при $a\ll\lambda_D\ll\lambda_{e,i}$~\cite[с.~9, (1)]{Romashchenko2021}.

\begin{outofcorpus}[OML у магнітному полі]
У полі токамака $\rho_e\sim$~мкм, тобто порівнянний з розміром пилинки, тож OML для електронів~--- лише наближення.
\end{outofcorpus}

\begin{corpusexample}[DUSTT/UEDGE: DIII-D і умови росту~\cite{Pigarov2006}]
У DUSTT/UEDGE (DIII-D) мікронні пилинки C із зовнішньої мішені розганяє тертя з іонами (у середньому \textbf{400~м/с}); у гарячій плазмі вони сублімують, заряджаються слабо додатно (емісія) й можуть сягати ядра~\cite[с.~5--6]{Pigarov2006}. Ріст обмежує температура пилинки: при $700<T_{d,\mathrm{eq}}<\SI{1800}{\degreeCelsius}$ хімічне розпилення вже пригнічене, а сублімація ще слабка. Для $T_e=\SI{10}{еВ}$, $n_e=\SI{e18}{м^{-3}}$ і 1\,\% C$^+$ маємо $T_{d,\mathrm{eq}}\approx\SI{1100}{\degreeCelsius}$ і ріст $\approx\SI{10}{нм/с}$; за кількох відсотків домішок~--- до \SI{100}{нм/с}~\cite[с.~6--7]{Pigarov2006}. Абляція пилу вже при конверсії осаду в пил $\fconv\approx0{,}01\,\%$ дає на сепаратрисі стільки ж нейтрального C, скільки розпилення~\cite[с.~9]{Pigarov2006}.
\end{corpusexample}

\section{Кондиціонування стінок тліючим і НВЧ-розрядом}\label{sec:08-gdc}

Між імпульсами стінки чистять тліючим розрядом (glow discharge, GDC), де катод~--- сама камера. Пробій: умова Таунсенда $\kT d=\ln(1+1/\gamma_{se})$, де $\kT$~[м$^{-1}$]~--- перший коефіцієнт Таунсенда (у літературі $\alpha$), $\kT/p=A_{\mathrm P}\ee^{-B_{\mathrm P}p/E}$, тобто крива Пашена~\cite[с.~54, (2.1), (2.3)]{Martsenyuk2025}; на Ураган-2М (калотні аноди) мінімум $\approx250$~В (Ar), 514~В (He), 380~В (H$_2$) при $\sim\SI{13}{Па}$~\cite[с.~52]{Martsenyuk2025}. 0D-модель Марценюка спирається на бомівський потік на стінку~\cite[с.~85, (3.1), (3.3)]{Martsenyuk2025}:
\begin{equation}\label{eq:08-global}
  \frac{v_B}{K_{iz}(T_e)}=\frac{N_0V}{A_{\mathrm{eff}}},\qquad
  n_e=\frac{P_{\mathrm{abs}}}{e\,v_BA_{\mathrm{eff}}(\mathcal E_{iz}+\mathcal E_e+\mathcal E_i)},\qquad
  \mathcal E_e=2T_e,\ \ \mathcal E_i=\tfrac12T_e+V_s .
\end{equation}
$v_B=\sqrt{eT_e/m_i}$, $A_{\mathrm{eff}}=A_Rh_R$; $T_e$ і всі $\mathcal E$ тут у вольтах. Баланс частинок задає $T_e$ незалежно від потужності, енергії~--- $n_e$; $2T_e$ і $\tfrac12T_e$~--- ті самі електронний потік енергії і розгін іонів у передшарі, що й у п.~\ref{sec:08-sheath}, $V_s\approx$ напруга розряду. При \SI{0,4}{Па} і \SI{162,5}{Вт} модель дає $T_e=\SI{2,19}{еВ}$ проти виміряних \SI{3}{еВ} (27\,\%) і $n_e=\SI{4e14}{м^{-3}}$ проти $\SI{3,7e14}{м^{-3}}$ (7,5\,\%)~\cite[с.~91]{Martsenyuk2025}. Додавання НВЧ (2,45~ГГц) знижує напругу: на Ураган-2М до 17--64~В ($\Delta U_d\approx145$--185~В) протягом $\sim\SI{1}{с}$, на TOMAS $\Delta U_d$ до $\approx\SI{220}{В}$ при $P_{\mathrm{mw}}\le\SI{1,5}{кВт}$ (майже лінійно в 0,4--1,2~кВт); менша напруга~--- менше розпилення самої стінки~\cite[с.~114]{Martsenyuk2025}.

\begin{practicum}[шар і двоточкова модель: \texttt{manual/code/ch08\_sheath\_2pm.py}]
\textbf{A.} Перевірте $\gamma_{\mathrm{sh}}(\mathrm D,T_i=T_e)=7{,}34$. \textbf{B.} Двоточкова модель ($Q_\parallel=\SI{30}{МВт/м^2}$, $L=\SI{40}{м}$, $\gamma_{\mathrm{sh}}=7$, ілюстративно): при $n_u=1\to\SI{4e19}{м^{-3}}$ $T_u\approx73\to64$~еВ, $T_t\approx56\to4{,}6$~еВ; перевірте нахили $-2$ і $+3$ (вони асимптотичні, при $n_u\gtrsim\SI{2e19}{м^{-3}}$). \textbf{C.} Додайте частку випромінювання $f_{\mathrm{rad}}$ і знайдіть $f_{\mathrm{rad}}$, що дає $T_t=\SI{5}{еВ}$ при $n_u=\SI{2e19}{м^{-3}}$; нижче \SI{5}{еВ} модель не працює (70\,\% енергії в надтеплових електронах~\cite{AhoMantila2011}).
\end{practicum}

\subsection*{Контрольні запитання}
\begin{enumerate}
  \item Чому падіння потенціалу в шарі пропорційне $T_e$ і лише логарифмічно залежить від маси іона? Як зміниться \eqref{eq:08-Vsh}, якщо коефіцієнт вторинної емісії $\delta_e\to1$?
  \item Чому в режимі високого рециклінгу $T_u$ майже не залежить від густини, а $T_t\propto n_u^{-2}$? Які припущення двоточкової моделі порушує детачмент і чому усереднена колізійність не є критерієм кінетичних ефектів~\cite{AhoMantila2011}?
  \item Чому W ерозують переважно нейтрали перезарядки й саморозпилення~\cite{AhoMantila2011,Belyaeva2016}? Чому W-rc кращий для дзеркал?
  \item Чому feedforward-подача Ar у DEMO надійніша за зворотний зв'язок~\cite{Fable2022}? Що спільного між бомівським потоком у 0D-моделі тліючого розряду~\cite{Martsenyuk2025} і граничною умовою SOLPS?
\end{enumerate}

\subsection*{Задачі}
\begin{enumerate}
  \item \emph{Шар.} D-плазма, $T_e=T_i=\SI{20}{еВ}$ (нижня межа режиму, обмеженого шаром~\cite{AhoMantila2011}). Обчисліть $eV_{\mathrm{sh}}$, повне падіння з передшаром, $\gamma_{\mathrm{sh}}$ і $E_{\mathrm{imp}}$ для D$^+$ та W$^{4+}$; порівняйте з порогом розпилення W воднем~\cite{AhoMantila2011}.
  \item \emph{Двоточкова модель.} $Q_\parallel=\SI{30}{МВт/м^2}$, $L=\SI{40}{м}$, $k_{0e}=2000$, $\gamma_{\mathrm{sh}}=7$, D. Знайдіть $T_u$, $T_t$, $n_t$ для $n_u=2$ і $\SI{4e19}{м^{-3}}$. У скільки разів треба збільшити напуск для подвоєння $n_u$ за~\cite{Fable2022}?
  \item \emph{Пилинка.} Вуглецева пилинка з $a=\SI{1}{мкм}$ у D-плазмі з $T_e=T_i=\SI{10}{еВ}$, $n=\SI{e18}{м^{-3}}$~\cite{Pigarov2006}. Перевірте умову OML~\cite{Romashchenko2021}, знайдіть плаваючий потенціал (без емісії), заряд за ємнісною моделлю і час подвоєння радіуса при рості \SI{10}{нм/с}.
\end{enumerate}
